% page 138 of the facsimile (image p-137 of the scan)
% block 165.1 x 242.0 mm ; left margin 8.0 mm
\pgc{-12.700mm}{0.000mm}{
\l{\cel{6}128}
\l{}
\l{}
\l{eklipsar.- (\sou{trans.)}\cel{2}I. (\sou{pri astro)}. Cesigar lua videbleso. -}
\l{\cel{3}II. (\sou{metaf.)}\cel{1}Obskurigar per brilar supera-grade. - EFIRS.}
\l{}
\l{ekliptiko.- (\sou{astron.}) Orbito quan Tero parkuras cirkum Suno.}
\l{\cel{3}- DEFIRS.}
\l{}
\l{\sou{eklogo}. Mikra poemo en qua encenigesas muton-pastori.-DEFIRS.}
\l{}
\l{ekologio.- Cienco di la relati da la organismi kun la medio}
\l{\cel{3}en qua li vivas, kun la cirkonstanci organika (biologio)}
\l{\cel{3}e neorganika (kosmala) di la existo.}
\l{}
\l{\sou{ekonomo}. - En granda dom-edo od en institucuro : la persono di}
\l{\cel{3}qua la rolo konsistas ye jerar la spensi. - DFIRS.}
\l{}
\l{\sou{ekonomiko}.- La cienco pri la produktado, la cirkulado, la kam-}
\l{\cel{3}biado e la repartiso di la kozi valoroza. - DEFIRS.}
\l{}
\l{\sou{ekonomio}.- I. Administro financala di domo, societo o stato.}
\l{\cel{3}- II. (\sou{metaf}.) (\sou{pri organismo}). Gani (nutrado) e spensi (per-}
\l{\cel{3}di di materio e di energio). - DEFIRS.}
\l{}
\l{\sou{ekritar}. - (\sou{trans}).(\sou{komerco}). Notar en la konto-libri la fluktui}
\l{\cel{3}di la aktivo e di la pasivo di omna konto. - F.}
\l{}
\l{\sou{ektodermo}.- (\sou{embriol}.). Strato extera di la blastodermo, qua pose}
\l{\cel{3}divenos la epidermo e la nervaro.}
\l{}
\l{\sou{ektoplasmo}.- I. (\sou{biol.}) Parto periferia di la celulo, diferencia-}
\l{\cel{3}jo de la citoplasmo, tale ke la kambii da ica kun la medio}
\l{\cel{3}eventas segun la naturo di la celulo. - II. (\sou{cienci okulta})}
\l{\cel{3}Plasmo, da origino psikala, qua (segun-dice) emanas de me-}
\l{\cel{3}diumo.}
\l{}
\l{\sou{ekumenika}.- (\sou{religio})\cel{2}Universala. - DEFIS.}
\l{}
\l{\sou{ekzemo.(patol.}) Morbo di la pelo, karakterizata da la erupto}
\l{\cel{3}di mikra vezikuli, sequata de exkorii surfacala. - DEFIRS.}
\l{}
\l{\sou{el(u)}.- Pronomo qua remplasas personino kande onu volas emfazar}
\l{\cel{3}ke lu esas femina.}
\l{}
\l{\sou{elaborar}. (\sou{trans.}). I.Transformar, produktar per laborado.}
\l{\cel{3}(\sou{anke metaf}.) II. (\sou{Fiziol}.) Igar absorbebla od asimilebla,}
\l{\cel{3}substanco od igar ica tala ke lu esas ekpulsebla de la}
\l{\cel{3}organismo o de la korpo. - EF.}
\l{}
\l{\sou{elastika}.- Qua esas kontraktebla, extensebla, deformebla per}
\l{\cel{3}preso, tiro, e c. e retrohavas sua volumino, sua formo ko-}
\l{\cel{3}mencala kande ta efiko cesas. - DEFIRS.}
\l{}
\l{\sou{elastino}.- Materio ek la klaso \textquotedbl{}albuminoidi\textquotedbl{}, quan onu ren-}
\l{\cel{3}kontras en la tegumenti ed en la subten-organi.}
}
